\subsection{Bilinear Pairings and Radicals}\label{subsec:metric-vector-spaces-and-bilinear-forms}

% **Role:** arbitrary-field metric geometry.
%
% Symmetric, alternating, and skew-symmetric bilinear forms; radicals; hyperbolic planes;
% symplectic and orthogonal geometries; Witt theory and canonical forms under congruence.

Metric vector spaces keep the algebraic part of inner-product geometry and drop positivity. Throughout this subsection $\mathbb{F}$ is an arbitrary field and vector spaces are finite-dimensional unless explicitly stated otherwise. The bilinear form is written $\innerprod{\mathbf{x},\mathbf{y}}$ when it is fixed.

\begin{definition} \label{def:bilinear-pairing-left-right-radicals} \hlt{Bilinear Pairings and Their Radicals}\\
Let $\mathcal{U},\mathcal{W}$ be vector spaces over $\mathbb{F}$.
\begin{enumerate}[label=(\roman*)]
\setlength{\itemsep}{0pt}
\item A \hlt{bilinear pairing} is a map $B:\mathcal{U}\times\mathcal{W}\to\mathbb{F}$ linear in each variable.
\item The \hlt{left radical} is $\operatorname{rad}_L(B)\defeq\{\mathbf{u}\in\mathcal{U}:B(\mathbf{u},\mathbf{w})=0\ \forall\,\mathbf{w}\in\mathcal{W}\}$.\\
The \hlt{right radical} is $\operatorname{rad}_R(B)\defeq\{\mathbf{w}\in\mathcal{W}:B(\mathbf{u},\mathbf{w})=0\ \forall\,\mathbf{u}\in\mathcal{U}\}$.
\item The \hlt{right orthogonal} of $S\subseteq\mathcal{U}$ is $S^{\perp_R}\defeq\{\mathbf{w}\in\mathcal{W}:B(\mathbf{s},\mathbf{w})=0\ \forall\,\mathbf{s}\in S\}$.\\
The \hlt{left orthogonal} of $T\subseteq\mathcal{W}$ is ${}^{\perp_L}T\defeq\{\mathbf{u}\in\mathcal{U}:B(\mathbf{u},\mathbf{t})=0\ \forall\,\mathbf{t}\in T\}$.
\item $B$ is \hlt{left nonsingular} if $\operatorname{rad}_L(B)=\{\mathbf{0}\}$, \hlt{right nonsingular} if $\operatorname{rad}_R(B)=\{\mathbf{0}\}$, and \hlt{perfect} if both hold.
\end{enumerate}
\end{definition}

\begin{proposition} \label{prop:bilinear-pairing-matrix-and-radicals} \hlt{Matrices of Bilinear Pairings}\\
Let $\alpha=(\mathbf{u}_1,\ldots,\mathbf{u}_m)$, $\beta=(\mathbf{w}_1,\ldots,\mathbf{w}_n)$ be bases of $\mathcal{U}$, $\mathcal{W}$. Let $M=M_{\alpha,\beta}\defeq [B(\mathbf{u}_i,\mathbf{w}_j)]_{i,j}\in\operatorname{Mat}(m,n,\mathbb{F})$.
\begin{enumerate}[label=(\roman*)]
\setlength{\itemsep}{0pt}
\item $B(\mathbf{u},\mathbf{w})=[\mathbf{u}]_{\alpha}^{t}M[\mathbf{w}]_{\beta}$.
\item $\{[\mathbf{u}]_{\alpha}:\mathbf{u}\in\operatorname{rad}_L(B)\}=\ker(M^t)$ and $\{[\mathbf{w}]_{\beta}:\mathbf{w}\in\operatorname{rad}_R(B)\}=\ker(M)$.
\item $\dim\operatorname{rad}_L(B)=m-\operatorname{rank}(M)$ and $\dim\operatorname{rad}_R(B)=n-\operatorname{rank}(M)$.
\item If $B$ is perfect, then $m=n$; when $m=n$, $B$ is perfect if and only if $M$ is invertible.
\end{enumerate}
\begin{proof}
\begin{enumerate}[label=(\roman*)]
\setlength{\itemsep}{0pt}
\item If $\mathbf{u}=\sum_i x_i\mathbf{u}_i$ and $\mathbf{w}=\sum_j y_j\mathbf{w}_j$, then $B(\mathbf{u},\mathbf{w})=\sum_{i,j}x_iM_{ij}y_j=[\mathbf{u}]_{\alpha}^{t}M[\mathbf{w}]_{\beta}$.
\item For $\mathbf{x}=[\mathbf{u}]_{\alpha}$ and $\mathbf{y}=[\mathbf{w}]_{\beta}$, $\mathbf{u}\in\operatorname{rad}_L(B)$ $\Longleftrightarrow$ $\mathbf{x}^tM\mathbf{y}=0\ \forall\,\mathbf{y}$ $\Longleftrightarrow$ $M^t\mathbf{x}=\mathbf{0}$.\\
Similarly, $\mathbf{w}\in\operatorname{rad}_R(B)\Longleftrightarrow M\mathbf{y}=\mathbf{0}$.
\item By (ii) and rank-nullity, $\dim\operatorname{rad}_L(B)=\dim\ker(M^t)=m-\operatorname{rank}(M^t)=m-\operatorname{rank}(M)$,\\
$\dim\operatorname{rad}_R(B)=\dim\ker(M)=n-\operatorname{rank}(M)$.
\item By (iii), perfection is equivalent to $m-\operatorname{rank}(M)=0$ and $n-\operatorname{rank}(M)=0$. Hence $m=n$; when $m=n$, this is equivalent to $M$ being invertible.
\end{enumerate}
\hlt{Structural Commentary}: A pairing matrix has one radical for each input slot.\\
\hlt{Key Insight}: Nonsingularity of a pairing is ordinary full-rank behavior of its representing matrix.
\end{proof}
\end{proposition}

\begin{proposition} \label{prop:bilinear-pairing-change-of-bases} \hlt{Two-Sided Basis Change for Pairings}\\
With notation as above, let $\alpha'$ and $\beta'$ be bases, and set $P=C_{\alpha\leftarrow\alpha'}$, $Q=C_{\beta\leftarrow\beta'}$. Then $M_{\alpha',\beta'}=P^tM_{\alpha,\beta}Q$.\\
For square forms this is congruence; for rectangular pairings it is two-sided equivalence.
\begin{proof}
By Proposition \ref{prop:bilinear-pairing-matrix-and-radicals}, $[\mathbf{u}]_{\alpha'}^tM_{\alpha',\beta'}[\mathbf{w}]_{\beta'} =B(\mathbf{u},\mathbf{w}) =(P[\mathbf{u}]_{\alpha'})^tM_{\alpha,\beta}(Q[\mathbf{w}]_{\beta'}) =[\mathbf{u}]_{\alpha'}^tP^tM_{\alpha,\beta}Q[\mathbf{w}]_{\beta'}$.\\
Thus $M_{\alpha',\beta'}=P^tM_{\alpha,\beta}Q$.\\
\hlt{Structural Commentary}: Pairing matrices change independently in each input slot.\\
\hlt{Key Insight}: The number of input slots determines how basis changes act on the matrix.
\end{proof}
\end{proposition}

\begin{perspective} \label{per:metric-spaces-specialize-bilinear-pairings} \hlt{Square Pairings}\\
Henceforth $B:\mathcal{V}\times\mathcal{V}\to\mathbb{F}$ is a square pairing, usually symmetric, skew-symmetric, or alternating.\\
In these cases $\operatorname{rad}_L(B)=\operatorname{rad}_R(B)$, so one radical and one orthogonality notation suffice.
\end{perspective}

\begin{definition} \label{def:metric-vector-space-basic-form-types} \hlt{Symmetric, Skew-Symmetric, and Alternating Forms}\\
Let $\mathcal{V}$ be a vector space over $\mathbb{F}$ and let $B:\mathcal{V}\times\mathcal{V}\to\mathbb{F}$ be bilinear.
\begin{enumerate}[label=(\roman*)]
\setlength{\itemsep}{0pt}
\item $B$ is a \hlt{bilinear form} on $\mathcal{V}$; when fixed, write $B(\mathbf{x},\mathbf{y})=\innerprod{\mathbf{x},\mathbf{y}}$.
\item $\operatorname{Sym}(B)$ means $B(\mathbf{x},\mathbf{y})=B(\mathbf{y},\mathbf{x})$ for all $\mathbf{x},\mathbf{y}\in\mathcal{V}$; then $B$ is \hlt{symmetric}.
\item $\operatorname{Skew}(B)$ means $B(\mathbf{x},\mathbf{y})=-B(\mathbf{y},\mathbf{x})$ for all $\mathbf{x},\mathbf{y}\in\mathcal{V}$; then $B$ is \hlt{skew-symmetric}.
\item $\operatorname{Alt}(B)$ means $B(\mathbf{x},\mathbf{x})=0$ for all $\mathbf{x}\in\mathcal{V}$; then $B$ is \hlt{alternating}.
\item A \hlt{metric vector space} is a pair $(\mathcal{V},B)$ with $\operatorname{Sym}(B)$, $\operatorname{Skew}(B)$, or $\operatorname{Alt}(B)$.\\
The symmetric case is an \hlt{orthogonal geometry}; the alternating case is a \hlt{symplectic geometry}.
\end{enumerate}
\end{definition}

\begin{proposition} \label{prop:form-type-relations-by-characteristic} \hlt{Characteristic Controls the Form Types}\\
Let $B(\mathbf{x},\mathbf{y})=\innerprod{\mathbf{x},\mathbf{y}}$ be a bilinear form on $\mathcal{V}$.
\begin{enumerate}[label=(\roman*)]
\setlength{\itemsep}{0pt}
\item $\operatorname{Alt}(B)\Rightarrow\operatorname{Skew}(B)$.
\item If $\operatorname{char}(\mathbb{F})\neq2$, then $\operatorname{Alt}(B)\Longleftrightarrow\operatorname{Skew}(B)$, and $\operatorname{Sym}(B)\wedge\operatorname{Alt}(B)\Rightarrow B=0$.
\item If $\operatorname{char}(\mathbb{F})=2$, then $\operatorname{Sym}(B)\Longleftrightarrow\operatorname{Skew}(B)$ and $\operatorname{Alt}(B)\Rightarrow\operatorname{Sym}(B)$.
\end{enumerate}
\begin{proof}
\begin{enumerate}[label=(\roman*)]
\setlength{\itemsep}{0pt}
\item If $\operatorname{Alt}(B)$, then $0=\innerprod{\mathbf{x}+\mathbf{y},\mathbf{x}+\mathbf{y}} =\innerprod{\mathbf{x},\mathbf{y}}+\innerprod{\mathbf{y},\mathbf{x}}$, so $\innerprod{\mathbf{x},\mathbf{y}}=-\innerprod{\mathbf{y},\mathbf{x}}$.
\item Assume $\operatorname{char}(\mathbb{F})\neq2$. By (i), $\operatorname{Alt}(B)\Rightarrow\operatorname{Skew}(B)$. Conversely, $\operatorname{Skew}(B)$ gives $\innerprod{\mathbf{x},\mathbf{x}}=-\innerprod{\mathbf{x},\mathbf{x}}$, hence $2\innerprod{\mathbf{x},\mathbf{x}}=0$ and $\operatorname{Alt}(B)$. If also $\operatorname{Sym}(B)$ and $\operatorname{Alt}(B)$, then $\innerprod{\mathbf{x},\mathbf{y}}=-\innerprod{\mathbf{x},\mathbf{y}}$, so $\innerprod{\mathbf{x},\mathbf{y}}=0$ for all $\mathbf{x},\mathbf{y}$.
\item If $\operatorname{char}(\mathbb{F})=2$, then $-1=1$, so $\operatorname{Sym}(B)\Longleftrightarrow\operatorname{Skew}(B)$.\\
Moreover $\operatorname{Alt}(B)\Rightarrow\operatorname{Skew}(B)$ by (i), hence $\operatorname{Alt}(B)\Rightarrow\operatorname{Sym}(B)$.
\end{enumerate}
\hlt{Structural Commentary}: Skew-symmetry collapses by characteristic.\\
\hlt{Key Insight}: The equation $\innerprod{\mathbf{x}+\mathbf{y},\mathbf{x}+\mathbf{y}}=0$ polarizes alternation.
\end{proof}
\end{proposition}

\begin{definition} \label{def:minkowski-metric-vector-space} \hlt{Minkowski Space}\\
\hlt{Minkowski space} is the orthogonal geometry $\mathbb{R}^{3,1}\defeq\mathbb{R}^4$ with standard basis $(\mathbf{e}_1,\ldots,\mathbf{e}_4)$\\
and bilinear form $\innerprod{\mathbf{x},\mathbf{y}} =x_1y_1+x_2y_2+x_3y_3-x_4y_4$.\\
Equivalently, its Gram matrix is $\operatorname{diag}(1,1,1,-1)$. It is nonsingular but not positive definite.
\end{definition}

\begin{example} \label{ex:minkowski-null-cone-and-inertia} \hlt{Null Vectors in Minkowski Space}\\
In $\mathbb{R}^{3,1}$, set $q(\mathbf{x})\defeq\innerprod{\mathbf{x},\mathbf{x}} =x_1^2+x_2^2+x_3^2-x_4^2$.\\
Null cone is $\mathcal{N}\defeq\{\mathbf{x}\neq\mathbf{0}:q(\mathbf{x})=0\} =\{\mathbf{x}\neq\mathbf{0}:x_1^2+x_2^2+x_3^2=x_4^2\}$. Every $\mathbf{x}\in\mathcal{N}$ is isotropic; i.e., $\mathbf{e}_1+\mathbf{e}_4\in\mathcal{N}$.\\
As $\operatorname{diag}(1,1,1,-1)$, $I_4$ have different real inertias, $(3,1,0)$ and $(4,0,0)$, $\mathbb{R}^{3,1}$ is not isometric to Euclidean $\mathbb{R}^4$.
\end{example}

\begin{definition} \label{def:matrix-rank-and-discriminant-of-form} \hlt{Matrix, Rank, and Discriminant of a Form}\\
Let $\beta=(\mathbf{v}_1,\ldots,\mathbf{v}_n)$ be an ordered basis of a metric vector space $\mathcal{V}$.
\begin{enumerate}[label=(\roman*)]
\setlength{\itemsep}{0pt}
\item The \hlt{matrix of the form} in $\beta$ is $M_{\beta}\defeq[\innerprod{\mathbf{v}_i,\mathbf{v}_j}]_{i,j}$.
\item The \hlt{rank} of the form is $\operatorname{rank}(M_{\beta})$, independent of $\beta$.
\item The \hlt{discriminant set} is $\Delta(\mathcal{V})\defeq\{r^2\det(M_{\beta}):r\in\mathbb{F}^{\times}\}$.
\item If the form is nonsingular, then $\Delta(\mathcal{V})$ is the square class of $\det(M_\beta)$ in $\mathbb{F}^{\times}/(\mathbb{F}^{\times})^2$.
\end{enumerate}
\end{definition}

\begin{proposition} \label{prop:metric-form-matrix-recovery-and-congruence} \hlt{Matrices Represent Forms up to Congruence}\\
Let $\beta,\gamma$ be ordered bases of $\mathcal{V}$ and put $C=C_{\beta\leftarrow\gamma}$.
\begin{enumerate}[label=(\roman*)]
\setlength{\itemsep}{0pt}
\item $\innerprod{\mathbf{x},\mathbf{y}}=[\mathbf{x}]_{\beta}^{t}M_{\beta}[\mathbf{y}]_{\beta}$.
\item $M_{\gamma}=C^tM_{\beta}C$.
\item Matrices $A,N$ represent the same form on $\mathcal{V}$ if and only if $N=P^tAP$ for some invertible $P$.
\end{enumerate}
\begin{proof}
\begin{enumerate}[label=(\roman*)]
\setlength{\itemsep}{0pt}
\item If $\mathbf{x}=\sum\limits_i a_i\mathbf{v}_i$ and $\mathbf{y}=\sum\limits_j b_j\mathbf{v}_j$, then $\innerprod{\mathbf{x},\mathbf{y}} =\sum\limits_{i,j}a_i b_j\innerprod{\mathbf{v}_i,\mathbf{v}_j} =[\mathbf{x}]_{\beta}^{t}M_{\beta}[\mathbf{y}]_{\beta}$.
\item Since $[\mathbf{x}]_{\beta}=C[\mathbf{x}]_{\gamma}$ and $[\mathbf{y}]_{\beta}=C[\mathbf{y}]_{\gamma}$, $[\mathbf{x}]_{\gamma}^{t}M_{\gamma}[\mathbf{y}]_{\gamma} =\innerprod{\mathbf{x},\mathbf{y}} =[\mathbf{x}]_{\gamma}^{t}C^tM_{\beta}C[\mathbf{y}]_{\gamma}$, so $M_{\gamma}=C^tM_{\beta}C$.
\item If $A=M_{\beta}$ and $N=M_{\gamma}$ represent one form, then (ii) gives $N=C_{\beta\leftarrow\gamma}^tAC_{\beta\leftarrow\gamma}$.\\
Conversely, suppose $N=P^tAP$ with $P=[p_{ij}]$ invertible and $A=M_{\beta}$.\\
Define $\gamma=(\mathbf{w}_1,\ldots,\mathbf{w}_n)$ by $\mathbf{w}_j=\sum_i p_{ij}\mathbf{v}_i$. Then $C_{\beta\leftarrow\gamma}=P$, so (ii) gives $M_{\gamma}=P^tAP=N$.
\end{enumerate}
\hlt{Structural Commentary}: Congruence is basis change for forms.\\
\hlt{Key Insight}: Basis change acts in both slots: $M\mapsto P^tMP$.
\end{proof}
\end{proposition}

